% ============================================================
\section{Algorithmen \& Datenstrukturen (Labs)}
% ============================================================
Aus den Labs 7 und 9: einfache Sortierverfahren und Graphen. Erg\"anzt die STL-Algorithmen aus dem Abschnitt \emph{Generische Programmierung \& STL}.

\subsection{Sortierverfahren}
Beide Verfahren sortieren \textbf{in-place} und haben die Zeitkomplexit\"at $O(n^2)$. Sie eignen sich zum Verstehen der Grundidee -- in der Praxis nutzt man \code{std::sort} ($O(n\log n)$).

\subsubsection*{Insertion Sort (Einf\"ugen)}
Jedes Element wird an die richtige Stelle des bereits sortierten linken Teils ,,einsortiert''. Best Case $O(n)$ (fast sortiert), \textbf{stabil}.
\begin{lstlisting}
for (std::size_t i = 1; i < values.size(); ++i) {
    int key = values[i];
    std::size_t j = i;
    while (j > 0 && values[j - 1] > key) {   // groessere Elemente nach rechts schieben
        values[j] = values[j - 1];
        --j;
    }
    values[j] = key;                          // key an richtiger Stelle einfuegen
}
\end{lstlisting}

\subsubsection*{Selection Sort (Auswahl)}
Sucht in jedem Durchlauf das Minimum des Rests und tauscht es nach vorne. Immer $O(n^2)$ (auch im Best Case), \textbf{nicht stabil}.
\begin{lstlisting}
for (std::size_t i = 0; i + 1 < values.size(); ++i) {
    std::size_t min_index = i;
    for (std::size_t j = i + 1; j < values.size(); ++j)
        if (values[j] < values[min_index]) min_index = j;
    if (min_index != i)
        std::swap(values[i], values[min_index]);   // Minimum nach vorne
}
\end{lstlisting}
\begin{center}\small
\begin{tabular}{@{}llll@{}}
\toprule
\textbf{Verfahren} & \textbf{Best} & \textbf{Average/Worst} & \textbf{Stabil?} \\
\midrule
Insertion Sort & $O(n)$ & $O(n^2)$ & ja \\
Selection Sort & $O(n^2)$ & $O(n^2)$ & nein \\
\code{std::sort} & \multicolumn{2}{c}{$O(n\log n)$} & nein \\
\bottomrule
\end{tabular}
\end{center}

\subsection{Graphen}
\begin{definition}[: Graph]
Ein Graph $G=(V,E)$ besteht aus \textbf{Knoten} (Vertices $V$) und \textbf{Kanten} (Edges $E$), die paarweise Beziehungen darstellen. Eigenschaften: \emph{gewichtet/ungewichtet}, \emph{gerichtet/ungerichtet}, \emph{zusammenh\"angend}, \emph{Grad} eines Knotens (Anzahl Kanten).
\end{definition}

\subsubsection*{Adjazenzmatrix vs. Adjazenzliste}
\begin{center}\small
\begin{tabular}{@{}lll@{}}
\toprule
 & \textbf{Adjazenzmatrix} & \textbf{Adjazenzliste} \\
\midrule
Struktur & $n\times n$-Matrix, \code{[i][j]} & pro Knoten Liste der Nachbarn \\
Kante pr\"ufen & $O(1)$ & $O(\text{Grad})$ \\
Speicher & $O(V^2)$ & $O(V + E)$ \\
gut bei & dichten Graphen & d\"unnen (sparse) Graphen \\
\bottomrule
\end{tabular}
\end{center}
\begin{lstlisting}
// Adjazenzmatrix (gewichtet) ueber 2D-Vektor:
std::vector<std::vector<int>> matrix(n, std::vector<int>(n, 0));
matrix[i][j] = gewicht;                    // Kante i->j mit Gewicht

// Adjazenzliste (gewichtet): pro Knoten Liste von (Nachbar, Gewicht)
std::vector<std::vector<std::pair<int,int>>> adj(n);
adj[i].push_back({j, gewicht});
\end{lstlisting}

\subsubsection*{Dijkstra -- k\"urzester Pfad}
Findet die k\"urzesten Pfade von einem Startknoten zu allen anderen in einem \textbf{gewichteten} Graphen. Idee: stets den unbesuchten Knoten mit der kleinsten bekannten Distanz w\"ahlen und dessen Nachbarn aktualisieren (\emph{Relaxation}).
\begin{merke}
\begin{itemize}
  \item \textbf{Effizienter} als alle Pfade zu berechnen, weil bereits gefundene k\"urzeste Distanzen wiederverwendet werden (gierige Auswahl) statt jede Kombination zu pr\"ufen.
  \item Ben\"otigt \textbf{nicht-negative} Kantengewichte -- bei negativen Gewichten kann ein als ,,fertig'' markierter Knoten nachtr\"aglich noch verbessert werden, was Dijkstra nicht ber\"ucksichtigt.
\end{itemize}
\end{merke}
\clearpage
